\section{The on-chain solver}
\label{sec:solver}

\Cref{thm:main} gives a unique root in a bracket whose endpoints straddle
it. This section runs the solver: integer bisection, in bounded gas. The
conditions that make the theorem apply on-chain are stated as explicit
obligations.

\subsection{Integer bisection}
\label{sec:solver:bisection}

The solver bisects $E$ on $\mathcal{B}$ at integer (WAD) precision. Each
step halves the bracket and keeps the half whose endpoints still straddle
the sign change. The bracket holds $\lvert\mathcal{B}\rvert = x_{\max}+1$
integers, so after $K = \lceil \log_2(\lvert\mathcal{B}\rvert) \rceil$
steps it has collapsed to a single integer $x^\star$, the boundary with
$E(x^\star) \geq 0 > E(x^\star + 1)$. The published rate is
$y^\star = x^\star / N$, rounded toward the senior tranche.

Bisection, not Newton, for three reasons. $F$ is non-smooth, its slope
jumping at branch boundaries (\cref{lem:F}), so a tangent-following
method has no derivative to follow, whereas bisection needs only that
$E$ be non-increasing. The iteration count $K$ is fixed at deploy time,
so gas is constant with no data-dependent exit. And the loop is a short
sweep over already-audited valuation code, with no invariant inversion
and a unique root.

\subsection{Two brackets: solve, then check expressibility}
\label{sec:solver:bracket}

The solver uses the \emph{conservation bracket} $\mathcal{B}$ of
\cref{thm:main} as its search domain. Its endpoints straddle the root by
construction, so the search always brackets a unique answer. Each probe
evaluates the frozen-$L$ valuation in closed form, not the pool's
invariant solver, so no probe re-triggers the E-CLP domain check. That
check runs once, when $L$ is computed at solve start from the live
balances.

The root is the economically correct rate, but the deployed pool can only
\emph{quote} rates inside its price band $[\alpha, \beta]$. With $\alpha,
\beta$ the lower and upper bounds on the senior price in quote, the
expressible range is
\begin{equation}
  x \;\in\; [N\alpha,\ N\beta].
\end{equation}
These bounds are per-solve: the configured band constrains the
rate-scaled price, so it is recentered at the cached rate $y_0$ (and if
the pool sorts the senior share as token1, the senior bounds are the
reciprocals of the pool's configured bounds, swapped). After solving,
the kernel checks $x^\star \in [N\alpha, N\beta]$. If the root
falls outside, the rate is real but unquotable: the sync reverts rather than clipping to a band edge and
publishing a rate the pool cannot honour. This is the only
market-condition revert in the solve, and it is a genuine economic
signal (a senior loss larger than the pool's band can represent), not a
solver fragility.

\subsection{Why the solve always converges}
\label{sec:solver:convergence}

\Cref{thm:main} is a statement about the function $E$. For it to govern
the on-chain loop, three implementation invariants must hold. Each is
cheap and is listed as an obligation in \cref{sec:solver:obligations}.

\begin{description}
  \item[Frozen inputs.] Every input to $E$ ($\STraw$, $N$, $f$, the
    invariant $L$, the cached rate $y_0$, and the checkpoint) is read
    once at solve start and never re-read. Bisection compares $E$ at
    different probes and assumes one fixed function. Any external reads
    inside $F$ (the yield-model view, the block timestamp) are constant
    across the bisection, since it runs in one block against a frozen
    checkpoint, so no probe sees a different function.
  \item[Live supply.] $N \geq 1$. The map divides by $N$, so the supply
    must be positive. A minimum senior position minted at market
    creation and never redeemable keeps $N$ bounded away from zero for
    the life of the market, so the division is always defined.
  \item[Total $E$.] Every probe returns a value rather than reverting.
    $E$ has two parts. $P_A$ maps every $x \in \mathcal{B}$ to a senior
    balance inside the E-CLP's computable domain
    (\S\ref{sec:solver:bracket}), so it never reverts. $F$ reverts only
    on a checkpoint that violates NAV conservation, and every checkpoint
    the accountant commits conserves by construction
    (\cref{sec:overview:waterfall}). So on the reachable state space $E$
    is defined at every probe in $\mathcal{B}$.
\end{description}

Under these three invariants, the solve terminates in exactly $K$ steps
and returns the unique root. The only remaining reverts are external to
the solver: a paused market, a stale or reverting price oracle, or the
surrounding pool operation hitting its own Balancer limits. None of these
is a failure of the solve, and in each case the cached rate is left
untouched, so a stale rate can never be read after a failed sync.

\subsection{Tolerance and rounding}
\label{sec:solver:rounding}

Two policy choices fix the solver's behaviour:
\begin{description}
  \item[Iteration count] is the $K$ of \cref{sec:solver:bisection},
    hardcoded per market. Each step of $E$ lies in $\{-1, 0\}$
    (\cref{sec:existence:mono}), so $K$ halvings collapse the bracket to
    the single integer boundary $x^\star$. No residual or
    iteration-count check is needed.
  \item[Rounding] sends the final $x^\star$, and the published
    $y^\star = x^\star/N$, toward the senior tranche. Together with the
    floor inside $P_A$ (\cref{eq:PA}) and the round-down of the
    invariant, every rounding step favours the protected tranche.
\end{description}

\subsection{Implementation obligations}
\label{sec:solver:obligations}

The correctness argument rests on the following, each verifiable per
market. The first group makes $E$ the right function, the second makes
the solve converge, the third keeps every step conservative.

\begin{description}
  \item[Senior independence.] The senior raw NAV is marked from sources
    that do not read the pool (its BPT value, its spot price, or the
    published rate), so the loop has exactly one fixed-point variable.
  \item[Valuation wiring.] The senior leg is lifted to WAD by its
    decimal factor $10^{18 - d_s}$ ($d_s$ its native decimals) and the
    published senior rate. The quote leg is lifted by $10^{18 - d_q}$
    and its rate $q$.
  \item[Frozen curve.] $L$ is computed once per solve from the balances
    scaled at $y_0$, rounded down (the underestimate), and held fixed across the
    bisection. $V_{\max}$ is the value at the upper corner of this same
    frozen curve.
  \item[Clamp.] The valuation uses the integral form
    $\widehat{V}(y) = \int_0^y \min(D, N)$ (\cref{sec:setup:PA}), with $N$
    read at solve start. In closed form this is $N y$ below the rate
    where demand falls to $N$, and $V$ shifted by a constant above it
    (the split is unique when demand is monotone, as for the Gyro
    E-CLP). Capping the share count inside the point value,
    $b_q + y\min(D,N)$, is a different function that is not $N$-Lipschitz
    and breaks monotonicity.
  \item[Frozen inputs, live supply, total $E$.] As in
    \cref{sec:solver:convergence}: read once, $N \geq 1$, conserving
    checkpoint.
  \item[Conservative rounding.] Senior-favouring at every step: the
    invariant down, the valuation floor, and the published rate toward
    senior.
\end{description}
